% Requires bibliography key Willwacher2026 and the established uniform
% leading-lattice, dyadic-measure, and arithmetic kernel-axis results.
\begin{theorem}[The normalized pro-two image]
\label{thm:normalized-protwo-free}
Put \(M=\ker\lambda\) and
\[
 R_M=\{g\in M:f_g=1\text{ in the free pro-two group}\},
 \qquad P=M/R_M.
\]
Then \(P\) is free pro-two on a countable family converging to
the identity. This family can be chosen arithmetically.

More precisely, choose the square-corrected arithmetic family
\(a_3=\sigma,a_5,a_7,\ldots\), whose leading reductions are the
free generators of
\(\mathfrak{grt}_1(\mathbf F_2)\), and for which
\(a_{j+2}\) is \([h_*,a_j]\) followed by inverse-square corrections,
where \(\lambda_2(h_*)=-3\).
If \(F_r\) is the free pro-two group on
\(z_3,\ldots,z_{2r+1}\), with transition maps killing the last
generator, then
\begin{equation}
 \varprojlim_r F_r \ \xrightarrow{\ \sim\ }\ P,
 \qquad z_j\longmapsto a_jR_M.                          \label{eq:normalized-protwo-1}
\end{equation}
In particular \(H^2_{\rm cts}(P,\mathbf F_2)=0\).
\end{theorem}

\begin{proof}
Write \(P_j=\operatorname{in}_j(a_j)\).
The integral leading-basis theorem gives an isomorphism
\begin{equation}
 \operatorname{Lie}_{\mathbf Z_2}(Z_3,Z_5,\ldots)
 \longrightarrow\bigoplus_n L_n,\qquad Z_j\longmapsto P_j,
                                                               \label{eq:normalized-protwo-2}
\end{equation}
where the letters have weights \(j\) and the target bracket is
Ihara's bracket.
For injectivity, divide a nonzero homogeneous polynomial by the
largest common power of two in its Lie-basis coefficients.
Its nonzero reduction remains nonzero after evaluation, by
\cite[Theorem~1.3]{Willwacher}.
Surjectivity is the integral leading-basis assertion.

Let \(\mathcal P=\varprojlim F_r\).
It has the following universal property: a sequence \(b_j\to1\)
in any pro-two group \(Q\) determines a unique continuous map
\(\mathcal P\to Q\) sending \(z_j\) to \(b_j\).
Indeed, in any finite quotient of \(Q\), all sufficiently late
\(b_j\) are one; hence the required map factors through a
finite-rank \(F_r\). These maps are compatible over the finite
quotients. The generators topologically generate \(\mathcal P\),
by compactness and their generation in every \(F_r\).

Use the weighted integral Magnus expansion \(z_j=1+Z_j\), with
\(\operatorname{wt}(Z_j)=j\). It is faithful.
For completeness, a finite-rank pro-two Magnus expansion is
already faithful modulo two: a nonidentity element survives in
a finite two-group \(Q\), and substituting \(Z_j\mapsto q_j-1\)
in \(\mathbf F_2[Q]\) detects it.
The augmentation ideal of this group algebra is nilpotent.
Inductively, a central element of order two gives a square-zero
ideal, and the quotient is the group algebra of a smaller
two-group. Thus the power-series substitution is legitimate.
Faithfulness in \(\mathcal P\) follows by projection to a
finite-rank \(F_r\).

For \(w\ne1\) in \(\mathcal P\), let \(Q_w\ne0\) be its first
nonzero weighted Magnus term, of weight \(n\).
It is an ordinary integral Lie polynomial.
In fact
\(\Delta(Z_j)=Z_j\otimes1+1\otimes Z_j+Z_j\otimes Z_j\);
the last term has larger weight, so the lowest term of a
group-like element is primitive for the ordinary coproduct.
Over \(\mathbf Q_2\), decomposition by word length and the
Dynkin operator give Lie membership.
The leading coefficients one in the Lyndon basis then give
integral Lie membership.

The images \(a_jR_M\) converge to one, since their parameters
start in degree \(j\). They therefore define the map \eqref{eq:normalized-protwo-1}.
Its image is closed and contains all integral leading-basis
corrections, so successive approximation and compactness make
it surjective.

To prove injectivity, act on the original Magnus algebra and put
\(A_j=\Phi_{a_j}-1\).
This operator raises degree by at least \(j\), and its component
of degree shift \(j\) is
\[
 D_{P_j}(S)=0,\qquad D_{P_j}(T)=[T,P_j].
\]
Substitution \(Z_j\mapsto A_j\) into the weighted completed
algebra is well-defined and agrees with the group
action. Its first possible nonzero term on \(w\) is
\[
 Q_w(D_{P_3},D_{P_5},\ldots)
    =D_{\,Q_w(P_3,P_5,\ldots)}.
\]
Here the equality follows from
\([D_A,D_B]=D_{\{A,B\}}\).
The polynomial inside the last derivation is nonzero by \eqref{eq:normalized-protwo-2}.
Its derivation is nonzero as well: a homogeneous polynomial
commuting with \(T\) is a pure power of \(T\), whereas these Lie
polynomials vanish when \(S=0\).
Thus \(w\) acts nontrivially, proving \eqref{eq:normalized-protwo-1}.

Finally, every extension of \(P\) by a finite two-group splits.
In the extension group choose an open subgroup meeting the
finite kernel trivially. It maps homeomorphically to an open
subgroup of \(P\).
The late generators have unique lifts there converging to one;
choose arbitrary lifts of the finitely many remaining
generators. The universal property of \eqref{eq:normalized-protwo-1} gives a section.
In particular every central extension by \(\mathbf F_2\)
splits, proving the cohomological assertion.
\end{proof}

\begin{corollary}[Frattini action and characters of the pro-two image]
\label{cor:protwo-image-invariant-characters}
For the family in Theorem~\ref{thm:normalized-protwo-free}, put
\(V=P/\overline{P^2[P,P]}\) and \(T=h_*-1\).
Then
\begin{equation}
 V\simeq\mathbf F_2[[T]],\qquad
 [a_{3+2r}]=T^r[a_3],\qquad h_*=1+T.                   \label{eq:normalized-protwo-3}
\end{equation}
The dyadic mass
\(\mu=\rho_{2,2}^I/2:P\to\mathbf Z_2\) is surjective.
For \(P_0=\ker\mu\),
\begin{equation}
 N/(N\cap R_M)\simeq D/R_D\simeq P_0,                  \label{eq:normalized-protwo-4}
\end{equation}
and
\begin{equation}
 H^1(P,\mathbf F_2)^G=\mathbf F_2(\mu\bmod2),
 \qquad
 H^1(P_0,\mathbf F_2)^G=\mathbf F_2\eta_3.             \label{eq:normalized-protwo-5}
\end{equation}
Thus every invariant quadratic character of \(M,N,\) or \(D\)
that factors through its full pro-two parameter image already
factors through the dyadic beta quotient.
\end{corollary}

\begin{proof}
The Frattini quotient of the inverse limit in \eqref{eq:normalized-protwo-1} is the product
of its generator copies of \(\mathbf F_2\).
Indeed its closed Frattini subgroup maps onto the Frattini
subgroup in every finite-rank quotient, and compactness
identifies it with the intersection of their inverse images.
The inverse-square corrections disappear in this quotient,
so
\([a_{j+2}]=(h_*-1)[a_j]\), proving \eqref{eq:normalized-protwo-3}.
An invariant continuous functional is consequently determined
by its value at \([a_3]\).
The mass modulo two is invariant and takes value one there,
proving the first assertion of \eqref{eq:normalized-protwo-5}.

The image of \(N\) is \(P_0\): for a lift \(m\) of a point of
\(P_0\), the number \(z=\rho_2^I(m)/2\) has zero dyadic
component. The pro-two image of \(\sigma^z\) is therefore one,
and \(m\sigma^{-z}\in N\) gives a lift.
The image of \(D\) is also \(P_0\).
Choose an arithmetic lift, take its dyadic-primary power, and
cancel its even dyadic odd-prime Kummer coordinates by the
arithmetic kernel axes. Its coordinate at two is already zero.
This proves \eqref{eq:normalized-protwo-4}.

Let \(\ell_r\) be the continuous coordinate functional dual to
\([a_{3+2r}]\). Then
\(H^1(P,\mathbf F_2)=\bigoplus_{r\ge0}\mathbf F_2\ell_r\),
with \(\ell_0=\mu\bmod2\).
The procyclic closure of \(a_3\) splits \(P\to\mathbf Z_2\).
Restriction thus gives the \(G\)-equivariant exact sequence
\[
 0\longrightarrow\mathbf F_2\ell_0
 \longrightarrow H^1(P,\mathbf F_2)
 \longrightarrow H^1(P_0,\mathbf F_2)^P
 \longrightarrow0.
\]
Any \(G\)-invariant character of \(P_0\) is \(P\)-invariant.
For \(\ell=\sum c_r\ell_r\), invariance of its class modulo
\(\ell_0\) under \(h_*\) requires
\[
 \ell\circ(h_*-1)=\sum_{r\ge1}c_r\ell_{r-1}
                   \in\mathbf F_2\ell_0.
\]
It follows that \(c_r=0\) for \(r\ge2\), so this invariant
space has dimension at most one.
The measure theorem gives the nonzero invariant character
\(\eta_3=m_3/8\bmod2\) on \(P_0\), already nonzero on
\(D\cap H\). This proves the remaining assertion of \eqref{eq:normalized-protwo-5}.
\end{proof}

\begin{remark}
The same finite-extension splitting argument gives the exact
sequence
\[
 0\longrightarrow H^1(P,\mathbf F_2)
 \longrightarrow H^1(M,\mathbf F_2)
 \longrightarrow H^1(R_M,\mathbf F_2)^M
 \longrightarrow0.
\]
After taking \(G/M\)-invariants, a possible obstruction remains
in \(H^1_{\rm cts}(G/M,H^1(P,\mathbf F_2))\).
These results do not prove that the relative abelian quotient
of \(R_D\) is zero.
\end{remark}
